\chapter{Динамічний та поведінковий аналіз шкідливого коду}
\chaptershort{Динамічний та поведінковий аналіз}
\label{chap:ch07}

Динамічний аналіз --- дослідження зразка під час його виконання в контрольованому середовищі --- доповнює статичний аналіз (розділ~\ref{chap:ch06}) там, де той безсилий: проти пакування, шифрування та обфускації (розділ~\ref{chap:ch08}), що приховують справжній код на диску. Під час виконання код мусить розгорнути себе в пам'яті й виявити справжню поведінку --- саме цю поведінку й спостерігає динамічний аналіз. Його сила --- у тому, що він фіксує \emph{що} зразок робить, незалежно від того, \emph{як} цей код записаний і захований; його слабкість --- у тому, що спостерігається лише \emph{одна} гілка виконання за даних умов середовища, а не всі можливі (на відміну від статичного аналізу), і в тому, що сучасний ЗПК активно виявляє середовище спостереження й ухиляється від нього. Розділ розглядає принципи й місце динамічного аналізу, архітектуру пісочниць, спостережувані артефакти поведінки, аналіз пам'яті, формалізацію й кластеризацію поведінки, межі методу й протидію антианалізним прийомам, і завершується кейсом дослідження зразка-вимагача. Поведінкове виявлення, що ґрунтується на динамічних ознаках, є основою сучасних засобів захисту кінцевих точок (розділ~\ref{chap:ch10}).

\section{Принципи й місце динамічного аналізу}\label{sec:ch07-principles}

Ідея динамічного аналізу проста: замість того щоб розбирати код, його \emph{запускають} у середовищі, де кожну дію можна спостерігати й за потреби обернути. Це знімає проблему прихованого коду, бо розпакування й розшифрування відбуваються автоматично --- їх виконує сам зразок. Натомість з'являються нові виклики: безпека (виконання реального ЗПК потребує ізоляції), повнота (одне виконання розкриває лише одну гілку), вірогідність (середовище спостереження має бути нерозрізнюваним від реального, інакше зразок не виявить справжньої поведінки). Систематичний огляд методів і засобів динамічного аналізу~\cite{ch07:egele} виокремлює ці компроміси як визначальні для галузі.

\begin{tikzfig}
\begin{tikzpicture}[font=\scriptsize,
  b/.style={draw,rounded corners=2pt,minimum height=6mm,text width=22mm,align=center,font=\tiny},
  >={Stealth[length=1.5mm]}]
\node[b] (in) at (0,0) {зразок};
\node[b] (st) at (0,-1.3) {статичний: усі гілки, прихований код};
\node[b] (dy) at (3.6,-1.3) {динамічний: одна гілка, справжня поведінка};
\node[b] (out) at (1.8,-2.8) {об'єднаний профіль загрози};
\draw[->] (in) -- (st);
\draw[->] (in) -- (dy);
\draw[->] (st) -- (out);
\draw[->] (dy) -- (out);
\draw[<->,dashed] (st) -- node[above,font=\tiny]{доповнення} (dy);
\end{tikzpicture}
\figcap{Рис.~\thechapter.1 --- Взаємодоповнення статичного й динамічного аналізу: кожен долає обмеження іншого}
\end{tikzfig}

Порівняння двох підходів наведено в табл.~\thechapter.1. Головний висновок: вони не конкурують, а доповнюють один одного, і зрілий аналіз є \term{гібридним} --- статичний спрямовує динамічний (визначає точки інтересу, умови запуску), а динамічний розгортає код для подальшого статичного аналізу дампа пам'яті.

\tabhead{Таблиця~\thechapter.1 --- Статичний та динамічний аналіз: порівняння}

\begin{center}\footnotesize
\begin{tabularx}{\linewidth}{L{22mm} Y Y}
\toprule
Властивість & Статичний & Динамічний \\
\midrule
Покриття гілок & усі гілки коду & лише виконані \\
Прихований код & вразливий до пакування & розгортається самим зразком \\
Безпека & код не виконується & потрібна ізоляція \\
Швидкість & висока & обмежена часом виконання \\
Ухилення & антидизасемблювання & антисендбокс, антиВМ \\
\bottomrule
\end{tabularx}
\end{center}

\section{Пісочниці та середовища спостереження}\label{sec:ch07-sandbox}

\term{Пісочниця} (sandbox) --- ізольоване інструментоване середовище для виконання й спостереження зразка. Історично першими автоматизованими пісочницями були системи на кшталт CWSandbox~\cite{ch07:willems} та Cuckoo~\cite{ch07:cuckoo}, що перехоплювали системні виклики й будували звіт про поведінку. Архітектурно пісочниці різняться за \emph{рівнем} інструментування та \emph{рівнем} ізоляції.

За рівнем спостереження розрізняють: \emph{перехоплення в користувацькому просторі} (хуки API-функцій; просте, але легко виявляється й обходиться прямими системними викликами); \emph{перехоплення в ядрі} (драйвери-монітори; стійкіше, але вразливе до руткітів рівня ядра, розділ~\ref{chap:ch05}); \emph{емуляцію та віртуалізацію з інтроспекцією} (аналіз ззовні гостьової системи через гіпервізор --- virtual machine introspection, що найважче виявити зсередини). За середовищем виконання --- емулятори, віртуальні машини та «голе залізо» (bare-metal): останнє застосовують саме проти зразків, які виявляють віртуалізацію~\cite{ch07:kirat}.

\begin{tikzfig}
\begin{tikzpicture}[font=\scriptsize,
  l/.style={draw,rounded corners=2pt,minimum height=5.5mm,text width=46mm,align=center,font=\tiny},
  >={Stealth[length=1.4mm]}]
\node[l] (s) at (0,0) {зразок (гостьова ОС)};
\node[l] (h) at (0,-1.0) {перехоплення: API / ядро / гіпервізор};
\node[l] (log) at (0,-2.0) {журнал подій (виклики, файли, реєстр, мережа)};
\node[l] (iso) at (0,-3.0) {ізоляція: ВМ / емулятор / bare-metal};
\node[l] (rep) at (0,-4.0) {звіт: поведінковий профіль, IoC};
\draw[->] (s) -- (h); \draw[->] (h) -- (log); \draw[->] (log) -- (rep);
\draw[->] (iso.west) to[out=180,in=180] (s.west);
\end{tikzpicture}
\figcap{Рис.~\thechapter.2 --- Узагальнена архітектура пісочниці для динамічного аналізу}
\end{tikzfig}

Якість пісочниці визначається двома суперечливими вимогами: \emph{прозорістю} (невидимістю для зразка) та \emph{повнотою спостереження}. Що глибше інструментування, то повніший звіт, але й то помітніше втручання. Що краще ізоляція імітує реальну систему (встановлені програми, історія діяльності користувача, реалістична мережа), то ймовірніше зразок виявить справжню поведінку.

Окремо слід розрізняти \emph{інтерактивний} аналіз, який веде аналітик за допомогою зневаджувача (debugger), покроково контролюючи виконання й оглядаючи стан, та \emph{автоматизований} аналіз у пісочниці, що масштабується на тисячі зразків за добу без участі людини. Перший незамінний для глибокого дослідження пріоритетних зразків і подолання складних прийомів антианалізу; другий --- основа промислового триажу потоку загроз (розділ~\ref{chap:ch03}). На практиці вони утворюють каскад, аналогічний ескалації глибини статичного аналізу (розділ~\ref{chap:ch06}): дешевий автоматизований прохід відсіює масу й визначає пріоритети, а дорогий ручний аналіз застосовують до того, що пройшло фільтр або чинить опір автоматиці.

Важливою архітектурною проблемою пісочниці є \emph{керування мережею}. Повна ізоляція без мережі позбавляє спостереження багатьох дій (звернення до командного сервера, завантаження додаткових модулів), бо зразок, не отримавши відповіді, завершується. Повний доступ до інтернету небезпечний (зразок може атакувати сторонні системи, а оператори --- виявити аналіз за зверненням з відомої аналітичної мережі й вимкнути сервер). Тому застосовують імітацію мережевих служб (підроблені відповіді DNS, HTTP, служб командного керування), що дозволяє спостерегти наміри зразка, не випускаючи трафік назовні, --- компроміс між повнотою й безпекою.

\section{Спостережувані артефакти поведінки}\label{sec:ch07-artifacts}

Поведінка зразка проявляється у взаємодії з операційною системою та мережею. Основні категорії спостережуваних артефактів: \emph{виклики API та системні виклики} (найбагатше джерело --- відображають наміри: робота з файлами, процесами, пам'яттю, реєстром, мережею, криптографією); \emph{файлова активність} (створення, зміна, видалення, шифрування файлів --- ключова ознака вимагачів); \emph{зміни реєстру} (персистентність, конфігурація --- розділ~\ref{chap:ch05}); \emph{мережева активність} (звернення до командних серверів, вивантаження даних, завантаження додаткових модулів); \emph{процесна активність} (породження процесів, ін'єкції, маніпуляції з токенами). Зіставлення цих артефактів із тактиками й техніками MITRE ATT\&CK~\cite{mitre-attack} дає структурований поведінковий профіль.

Особливу цінність має \emph{послідовність} і \emph{контекст} дій, а не окремі виклики. Один виклик шифрувальної функції нічого не означає; послідовність «перелік файлів $\to$ читання $\to$ шифрування $\to$ перезапис $\to$ видалення тіньових копій $\to$ створення записки про викуп» однозначно ідентифікує вимагача. Тому поведінкові сигнатури описують не окремі дії, а їх патерни --- подібно до того, як правила YARA (розділ~\ref{chap:ch06}) описують статичні ознаки.

Виклики API та системні виклики заслуговують окремої уваги як найбагатше й найстійкіше джерело. Будь-яка значуща дія --- робота з файлом, мережеве з'єднання, запис у реєстр, виділення виконуваної пам'яті, створення процесу --- рано чи пізно звертається до операційної системи через чітко визначений інтерфейс, який неможливо обійти, не відмовившись від самої дії. Це робить рівень системних викликів природним «вузьким місцем» для спостереження: хоч би як був обфускований код, його взаємодія зі світом проходить тут. Саме тому автори ЗПК застосовують прямі системні виклики в обхід перехоплених бібліотечних функцій, а розробники пісочниць переносять перехоплення глибше --- у ядро й гіпервізор, --- щоб це вузьке місце лишалося під контролем.

Мережева активність має подвійну цінність. По-перше, вона розкриває інфраструктуру зловмисника --- адреси й домени командних серверів, протоколи зв'язку, формати маячкових повідомлень (beacon), --- що безпосередньо живить мережеві індикатори компрометації (розділ~\ref{chap:ch12}) та виявлення на рівні мережі (розділ~\ref{chap:ch10}). По-друге, патерни трафіку (регулярність маячків, обсяг вивантажених даних, характерні заголовки) самі є поведінковою ознакою навіть за шифрованого вмісту. Тому мережеве спостереження --- невіддільна частина динамічного аналізу, а не окрема дисципліна.

\section{Аналіз пам'яті}\label{sec:ch07-memory}

\term{Аналіз пам'яті} (memory forensics) досліджує знімок оперативної пам'яті системи й посідає особливе місце: пам'ять містить \emph{розгорнутий} код (після розпакування), розшифровані дані, ключі, ін'єктовані модулі, мережеві з'єднання, сліди безфайлового коду, яких немає на диску. Фреймворк Volatility~\cite{ch07:volatility,ch07:ligh} став де-факто стандартом: він відновлює списки процесів (зокрема приховані руткітами через порівняння різних джерел --- cross-view, розділ~\ref{chap:ch05}), завантажені модулі, мережеві з'єднання, ін'єктовані ділянки пам'яті (з аномальними правами виконання), дескриптори.

Аналіз пам'яті критичний для виявлення \emph{безфайлового} коду (розділ~\ref{chap:ch11}), що існує лише в пам'яті й не залишає файлу на диску, а отже невидимий для статичного аналізу файлів. Він також дозволяє перемогти пакування пасивно: замість розпакування вручну аналітик дозволяє зразку розпакуватися самому, після чого знімає дамп і аналізує вже відкритий код статичними методами. Це робить аналіз пам'яті містком між динамічним і статичним аналізом.

Методологічно аналіз пам'яті спирається на реконструкцію структур ядра операційної системи зі «сирого» знімка: списків процесів і потоків, таблиць об'єктів, структур мережевого стеку, кешів файлової системи. Порівняння кількох незалежних джерел тієї самої інформації (cross-view, розділ~\ref{chap:ch05}) --- наприклад, списку процесів зі службового переліку диспетчера проти обходу структур ядра напряму --- викриває приховування: об'єкт, відсутній в одному поданні, але наявний в іншому, майже напевно прихований руткітом. Так аналіз пам'яті перетворює саме зусилля ЗПК зі приховування на детектувальну ознаку. Обмеження методу --- миттєвість знімка (він фіксує лише один момент і може не захопити короткочасні події) та потреба в надійному, незіпсованому самим ЗПК механізмі зняття дампа, бо руткіт рівня ядра здатен спотворити знімок, зроблений зсередини системи.

\section{Формалізація та кластеризація поведінки}\label{sec:ch07-formalize}

Для автоматизованого порівняння поведінку формалізують. Поширене подання --- \emph{множина} або \emph{послідовність} поведінкових ознак (викликів API, системних викликів, дій над ресурсами). Подібність поведінки двох зразків $A$ і $B$, поданих множинами ознак, оцінюють коефіцієнтом Жаккара~\eqref{eq:ch03-jaccard}, а для послідовностей --- мірами на $n$-грамах викликів. Поведінкова кластеризація~\cite{ch07:bayer,ch07:bailey} групує зразки за схожістю динамічних профілів, що стійкіше до поверхневих статичних змін, бо спирається на семантику дій.

Важливою характеристикою є \term{покриття} виконання --- частка коду (чи гілок), досягнута під час спостереження. Якщо позначити множину можливих семантичних гілок $\mathcal{B}$, а досягнуту за один запуск --- $\mathcal{B}_{\text{run}}$, то покриття

\begin{equation}
\mathrm{Cov} = \frac{|\mathcal{B}_{\text{run}}|}{|\mathcal{B}|} \in [0;1]
\label{eq:ch07-coverage}
\end{equation}

фундаментально обмежене: умовний код (зокрема логічні бомби, що спрацьовують за датою чи командою) та залежність від середовища залишають частину поведінки неспостереженою. Для підвищення покриття застосовують багаторазові запуски за різних умов та техніки примусового дослідження кількох гілок (multi-path execution).

Поведінкові профілі, отримані динамічно, мають і самостійну цінність для автоматичної класифікації. На відміну від статичних ознак, вони стійкі до пакування й поверхневих змін коду, бо відображають \emph{дії}, а не запис. Два зразки одного сімейства, скомпільовані по-різному чи запаковані різними пакувальниками (розділ~\ref{chap:ch08}), статично виглядають несхоже, але динамічно демонструють ту саму послідовність дій --- ту саму роботу з файлами, реєстром, мережею. Саме тому поведінкова кластеризація~\cite{ch07:bayer} стала основою масштабного триажу: вона групує функційно подібні зразки попри статичні відмінності, автоматично відносячи нові зразки до відомих сімейств. Обмеження --- та сама неповнота: якщо в спостереженому запуску сімейство не проявило характерної поведінки (через ухилення чи невиконану гілку), профіль буде неповним, а класифікація --- ненадійною. Тому поведінкову кластеризацію, як і решту динамічних методів, посилюють багаторазовим спостереженням і поєднанням зі статичними ознаками.

Ефективність спроб ухилення зручно описувати через імовірність того, що зразок \emph{не} виявить середовище спостереження. Якщо зразок застосовує $k$ незалежних перевірок, кожна з яких викриває пісочницю з імовірністю $p_i$, то ймовірність лишитися непоміченим і, отже, спостерегти справжню поведінку:

\begin{equation}
P_{\text{спост}} = \prod_{i=1}^{k} (1 - p_i).
\label{eq:ch07-evasion}
\end{equation}

Формула~\eqref{eq:ch07-evasion} пояснює стратегію вдосконалення пісочниць: зменшення кожного $p_i$ (підвищення прозорості окремих аспектів) та, у разі стійких до ВМ зразків, перехід на bare-metal, де низка перевірок віртуалізації дає $p_i \to 0$.

\section{Межі методу та протидія ухиленню}\label{sec:ch07-limits}

Головна вада динамічного аналізу --- неповнота: спостерігається лише виконана гілка за конкретних умов. Логічна бомба, що чекає певної дати чи команди, у пісочниці не спрацює й буде визнана нешкідливою. Друга проблема --- \term{антисендбокс} та \term{антиВМ} прийоми~\cite{ch07:chen,ch07:lindorfer}: зразок виявляє ознаки середовища спостереження (артефакти віртуалізації, малий час роботи системи, відсутність активності користувача, характерні імена пристроїв, присутність інструментів аналізу) і, виявивши їх, приховує справжню поведінку --- завершується, засинає або виконує безпечний код. Виявлення чутливих до середовища зразків~\cite{ch07:lindorfer} ґрунтується на порівнянні поведінки в різних середовищах: розбіжність сама є ознакою ухилення.

Арсенал антисендбокс-прийомів зручно поділити на кілька груп. \emph{Виявлення артефактів віртуалізації}: характерні рядки в описах пристроїв, специфічні інструкції процесора й розбіжності в їх поведінці, присутність драйверів і служб гіпервізора. \emph{Виявлення інструментування}: наявність перехоплень у функціях, процеси відомих аналітичних засобів, невідповідності в поверненнях API, спричинені хуками. \emph{Аналіз обстановки}: малий обсяг пам'яті чи кількість ядер, мало недавно відкритих документів і збережених облікових записів, «чиста» історія браузера, відсутність руху миші --- ознаки штучного середовища, а не реальної робочої станції. \emph{Часові й поведінкові пастки}: тривале присипляння в розрахунку на обмежений час аналізу, очікування перезавантаження чи певної дії користувача, спрацювання лише за конкретної дати (логічна бомба). Кожна група вимагає власних контрзаходів, і саме тому побудова прозорої пісочниці є не разовим рішенням, а безперервним процесом супроводу в міру появи нових перевірок.

Протидія ухиленню --- це постійна гонитва: підвищення прозорості пісочниць (інтроспекція через гіпервізор, реалістичні середовища з історією діяльності, емуляція взаємодії користувача), bare-metal-аналіз~\cite{ch07:kirat}, виявлення самих перевірок ухилення як підозрілої ознаки. Жодне середовище не є ідеально прозорим, тому динамічний аналіз, як і статичний, дає неповну картину й найефективніший у поєднанні з іншими методами. Вітчизняні дослідження систематизують застосування пісочниць~\cite{ch07:dynamic-ua} та поведінкових методів на основі системних викликів~\cite{ch07:behavior-ua} в умовах зростання кількості ухильних зразків.

\section{Кейс: динамічний аналіз зразка-вимагача}\label{sec:ch07-case}

Розглянемо послідовність динамічного аналізу підозрілого зразка (захисна задача: побудова поведінкового профілю для виявлення й реагування).

\emph{Крок 1. Підготовка середовища.} Готуємо ізольовану пісочницю з реалістичним профілем: встановлені офісні програми, документи-приманки, історія діяльності, контрольована мережа з імітацією інтернету (щоб спостерегти мережеві звернення, не випускаючи зразок назовні). Знімаємо базовий знімок системи.

\emph{Крок 2. Запуск і спостереження.} Запускаємо зразок, фіксуючи виклики API, файлову й реєстрову активність, мережеві звернення, породження процесів протягом заданого часу з емуляцією дій користувача.

\emph{Крок 3. Аналіз поведінки.} У журналі виявляємо характерну послідовність вимагача: перелік файлів користувача, читання й перезапис із високою ентропією вмісту (шифрування~\eqref{eq:ch03-entropy}), видалення тіньових копій, звернення до командного сервера по ключ, створення записок про викуп. Зіставляємо дії з техніками ATT\&CK.

\emph{Крок 4. Аналіз пам'яті.} Знімаємо дамп пам'яті й відновлюємо розшифрований код, рядки конфігурації, адресу командного сервера, можливий ключ --- артефакти, відсутні у файлі на диску.

\emph{Крок 5. Формування індикаторів і правил.} На основі профілю формуємо поведінкові правила виявлення (розділ~\ref{chap:ch11}), мережеві та файлові індикатори (розділ~\ref{chap:ch12}) для реагування (розділ~\ref{chap:ch16}). Перевіряємо стійкість: повторюємо запуск за інших умов (дати, мови системи), щоб виявити приховані гілки.

Кейс ілюструє силу динамічного аналізу --- розкриття справжньої поведінки попри пакування --- і його межі: висновок про нешкідливість за одним запуском ненадійний, тому критичні рішення спираються на багаторазове спостереження й поєднання зі статичним аналізом.

\section*{Висновки до розділу}
\addcontentsline{toc}{section}{Висновки до розділу}

Динамічний аналіз досліджує зразок під час виконання й долає головне обмеження статичного --- прихований пакуванням і шифруванням код, який зразок розгортає сам. Його сила --- у фіксації справжньої семантики дій незалежно від запису коду; слабкість --- у спостереженні лише однієї гілки за даних умов~\eqref{eq:ch07-coverage} та вразливості до антисендбокс-прийомів. Пісочниці балансують прозорість і повноту спостереження, розрізняючись рівнем перехоплення (API, ядро, гіпервізор) та ізоляції (емулятор, ВМ, bare-metal). Спостережувані артефакти --- виклики API, файлова, реєстрова, мережева й процесна активність --- найцінніші у вигляді контекстних \emph{послідовностей}, зіставлених із ATT\&CK. Аналіз пам'яті розкриває розгорнутий код, ключі та безфайлові загрози, слугуючи містком до статичного аналізу дампа. Формалізація поведінки уможливлює кластеризацію сімейств, стійкішу до поверхневих змін. Межі методу --- неповнота покриття й ухилення~\eqref{eq:ch07-evasion} --- долаються підвищенням прозорості, bare-metal-аналізом і багаторазовими запусками. Головний висновок: динамічний аналіз необхідний, але недостатній; максимальну цінність він дає в гібридному поєднанні зі статичним аналізом і є фундаментом поведінкового виявлення сучасних засобів захисту.

\selfcheck
\begin{enumerate}
\item У чому полягає принцип динамічного аналізу й яку головну проблему статичного аналізу він долає?
\item Порівняйте статичний і динамічний аналіз (табл.~\thechapter.1). Чому зрілий аналіз є гібридним?
\item Що таке пісочниця? Опишіть рівні перехоплення (API, ядро, гіпервізор) та їхні переваги й вади.
\item Поясніть суперечність між прозорістю й повнотою спостереження в пісочниці. Навіщо потрібне реалістичне середовище?
\item Які категорії артефактів поведінки спостерігають динамічно? Чому послідовність дій важливіша за окремі виклики?
\item Чому аналіз пам'яті посідає особливе місце? Що він розкриває з того, чого немає у файлі на диску?
\item Як аналіз пам'яті дозволяє перемогти пакування без ручного розпакування?
\item Що таке покриття виконання~\eqref{eq:ch07-coverage}? Чому воно принципово обмежене й як його підвищують?
\item За формулою~\eqref{eq:ch07-evasion} поясніть стратегію підвищення прозорості пісочниць. Що дає bare-metal-аналіз?
\item Що таке антисендбокс- та антиВМ-прийоми? Наведіть приклади ознак середовища, які виявляє зразок.
\item Як виявляють чутливі до середовища зразки? Чому розбіжність поведінки в різних середовищах є ознакою ухилення?
\item Чому висновок про нешкідливість зразка за одним запуском у пісочниці ненадійний?
\item Опишіть послідовність динамічного аналізу вимагача (кейс підрозділу~\ref{sec:ch07-case}). Яку роль відіграє аналіз пам'яті?
\item Дослідницьке: порівняйте повноту поведінкових профілів, отриманих за одного та за багаторазових запусків із варіацією умов, на наборі ухильних зразків.
\item Дослідницьке: оцініть стійкість обраної пісочниці до набору відомих антиВМ-перевірок і запропонуйте заходи підвищення її прозорості.
\end{enumerate}
